
\section{Decomposition theorem in the chiral setting} \label{section:chiral-setting}

\subsection{Affine W-algebra for the minimal orbit}

Affine W-algebras are defined by the following chiral Hamiltonian reduction process \cite{feigin1990quantization,kac2003quantum,kac2004quantum}. See also \cite[Section 4]{genra2025reduction} for a detailed survey of the construction. Let us define a chiral comoment map
\[ \widehat\mu \colon \V(\g_{\geqslant 1}) \longrightarrow \V^k(\g) \]
by mapping a generator $a$ in $\g_{\geqslant 1}$ to $a + \chi(a) 1$. It is a graded homomorphism because: $V(\g_{\geqslant 1}) \subseteq \V^k(\g)$ is a graded inclusion, and $\g_{\geqslant 1} = \g_1 \subseteq \V^k(\g)(0)$.

\begin{theorem}[{\cite[Theorem 4.1]{kac2004quantum}}]
	The cohomology of $\Cpx(\V^k(\g))$ is concentrated in degree $0$.
\end{theorem}

The degree $0$ of the cohomology is, by definition, the \emph{affine W-algebra} associated to $f$:
\[ \W^k(\g, f) \coloneqq \Hgy^0\big(\Cpx(\V^k(\g))\big). \]
According to \Cref{proposition:brst-cohomology}, the W-algebra is a graded vertex algebra whose grading is induced by \eqref{equation:grading-universal-affine}.

\begin{remark}
	Similarly to the finite W-algebra case, the fact that the grading is supported on $\Z_{\geqslant 0}$ comes from the fact that the good grading on $\g$ is supported on $\{-1, 0, 1\}$ (see \Cref{remark:nonnegative-filtration}).
\end{remark}

\subsection{The W-algebra is a vertex subalgebra of the localisation} \label{subsection:affine-w-algebra-coset}

We state a chiralisation of \Cref{theorem:slodowy-git}, that is the reconstruction of the affine W-algebra as a subalgebra of $\V^k(\g)[e^{-1}]$. 

The Lie algebra $\J_\infty(\kfk \ltimes \g_{\geqslant 1}) = (\kfk \ltimes \g_{\geqslant 1})[\![t]\!]$, as a subalgebra of $\J_\infty \g$, acts on $\V^k(\g)$ and the localisation $\V^k(\g)[e^{-1}]$. 

The Lie algebra $\J_\infty(\kfk \ltimes \g_{\geqslant 1})$ inherits a $\Z_{\leqslant 0}$-grading from the opposite good grading $\g = \bigoplus_{i \in \Z} \g_{-i}$ and by imposing that the formal variable $t$ is of degree $0$ (see \Cref{example:grading-jet-groups}). The properties of the Lie filtration (\Cref{theorem:li-filtration}) give that for $x$ in $\g_i \subseteq \V^k(\g)(1 - i)$, and for nonnegative integers $n$ and $p$,
\[ \begin{multlined}
	x t^n \cdot \big(\V^k(\g)[e^{-1}]\big)_p = x_{(n)} \big(\V^k(\g)[e^{-1}]\big)_p \\ \subseteq \big(\V^k(\g)[e^{-1}]\big)_{1 - i + p -1} = \big(\V^k(\g)[e^{-1}]\big)_{p - i}.
\end{multlined} \]
Therefore, $\V^k(\g)[e^{-1}]$ is a filtered $\J_\infty(\kfk \ltimes \g_{\geqslant 1})$-module in the sense of \Cref{subsection:filtered-representation}.

Because the action of $x t^n$ is given by the $n$-product $x_{(n)}$, we have:
\[ \begin{aligned}
	& \big(\V^k(\g)[e^{-1}] \big)^{\J_\infty(\kfk \ltimes \g_{\geqslant 1})} \\
	= & \big\{ F \in \V^k(\g)[e^{-1}] \quad | \quad \text{for all} ~  x \in \kfk \ltimes \g_{\geqslant 1} ~\text{and} ~n \geqslant 0,  ~  x_{(n)} F = 0 \big\} \\
	= & \big\{ F \in \V^k(\g)[e^{-1}] \quad | \quad \text{for all} ~  x \in \V^k(\kfk \ltimes \g_{\geqslant 1}),  ~   [x(z_1), F(z_2)] = 0 \big\}. 
\end{aligned} \]
Because it is defined as subalgebra commuting with a graded vertex subalgebra, $\big(\V^k(\g)[e^{-1}] \big)^{\J_\infty(\kfk \ltimes \g_{\geqslant 1})}$ is a graded vertex subalgebra of $\V^k(\g)[e^{-1}]$.

\begin{theorem} \label{theorem:affine-w-algebra-coset}
	There is a graded vertex algebra isomorphism
	\[ \big(\V^k(\g)[e^{-1}]\big)^{\J_\infty(\kfk \ltimes \g_{\geqslant 1})} \cong \W^k(\g, f) \]
\end{theorem}

To prove this theorem, we start by computing the associated graded algebra of $\big(\V^k(\g)[e^{-1}]\big)^{\J_\infty(\kfk \ltimes \g_{\geqslant 1})}$ for the filtration induced by the increasing Li filtration on $\V^k(\g)[e^{-1}]$.

\begin{proposition} \label{proposition:gr-invariant-borisov-localisation}
	The natural vertex Poisson algebra homomorphism
	\[ \gr\big((\V^k(\g)[e^{-1}])^{\J_\infty(\kfk \ltimes \g_{\geqslant 1})}\big) \longrightarrow \big(\gr(\V^k(\g)[e^{-1}])\big)^{\J_\infty(\kfk \ltimes \g_{\geqslant 1})} \]
	is an isomorphism. Therefore, we have a vertex Poisson algebra isomorphism:
	\[ \gr\big((\V^k(\g)[e^{-1}])^{\J_\infty(\kfk \ltimes \g_{\geqslant 1})}\big) \cong \J_\infty\C[\g^\vee[e^{-1}] \git (K \ltimes G_{\geqslant 1})].  \]
\end{proposition}

\begin{proof}
	For an algebraic group $N$, and for any $m$ in $\Z_{\geqslant 0}$, there is a semidirect product decomposition
	\[ \J_m N = N \ltimes (\J_m N)_1, \]
	where $(\J_m N)_1$ is the normal subgroup given by the kernel of the natural projection $\J_m N \rightarrow N$. The group $(\J_m N)_1$ is unipotent for finite $m$ and its Lie algebra is $t \n[t] / (t^{m+1})$.
	
	Therefore, the arc group $\J_\infty K = \J_\infty \Gm \times \J_\infty \Ga{}^{d-2}$ decomposes as the product of a torus and the projective limit of a system of unipotent groups:
	\[ \J_\infty K = \Gm \ltimes ((\J_\infty \Gm)_1 \times \J_\infty \Ga{}^{d-2}). \]
	
	To prove the proposition, we can first deal with the action of the normal subgroups: $((\J_\infty \Gm)_1 \times \J_\infty \Ga{}^{d-2}) \ltimes \J_\infty G_{\geqslant 1}$ in $\J_\infty K \ltimes \J_\infty G_{\geqslant 1}$, that is the projective limit of graded unipotent groups $((\J_m \Gm)_1 \times \J_m \Ga{}^{d-2}) \ltimes \J_m G_{\geqslant 1}$. We aim to show the isomorphism
	\begin{equation} \label{equation:gr-invariant-unipotent-jet} \gr \Big(\big(\V^k(\g)[e^{-1}]\big)^{(t \g_{\mathrm m}[\![t]\!] \oplus \g_{\mathrm a}{}^{d-2} [\![t]\!]) \ltimes \g_{\geqslant 1}[\![t]\!]}\Big) \cong \big(\gr \big(\V^k(\g)[e^{-1}]\big)\big)^{(t \g_{\mathrm m}[\![t]\!] \oplus \g_{\mathrm a}{}^{d-2} [\![t]\!]) \ltimes \g_{\geqslant 1}[\![t]\!]}. \end{equation}
	To do so, we use \Cref{proposition:gr-invariant-pro}. If we apply the arc functor $\J_\infty$ to the isomorphism of \Cref{proposition:dilation-gan-ginzburg}, we have the $((\J_\infty \Gm)_1 \times \J_\infty \Ga{}^{d-2}) \ltimes \J_\infty G_{\geqslant 1}$-equivariant isomorphism
	\begin{multline*} \C[((\J_\infty \Gm)_1 \times \J_\infty \Ga{}^{d-2}) \ltimes \J_\infty G_{\geqslant 1}] \otimes \C[\Gm \times \J_\infty S_f] \\ \cong \J_\infty (\Sym \g)[e^{-1}] \cong \gr (\V^k(\g)[e^{-1}]).  \end{multline*}
	The assumption of \Cref{proposition:gr-invariant-pro} about the Lie algebra gradings has been checked at the beginning of \Cref{subsection:affine-w-algebra-coset}. Therefore, we can apply the proposition and we get the isomorphism \eqref{equation:gr-invariant-unipotent-jet}.
	
	The isomorphism \eqref{equation:gr-invariant-unipotent-jet} is $\Gm$-equivariant, so we get the isomorphism
	\[ \Big(\big(\gr \big(\V^k(\g)[e^{-1}]\big)\big)^{(t \g_{\mathrm m}[t] \oplus \g_{\mathrm a}{}^{d-2} [t]) \ltimes \g_{\geqslant 1}[t]}\Big)^{\Gm} \cong \big(\gr \big(\V^k(\g)[e^{-1}]\big)\big)^{(\kfk \ltimes \g_{\geqslant 1})[t]} \]
	between the $\Gm$-invariant subalgebras. By applying \Cref{proposition:weight-filtration} on the left-hand side, we get the desired isomorphism:
	\[ \gr\big((\V^k(\g)[e^{-1}])^{\J_\infty(\kfk \ltimes \g_{\geqslant 1})}\big) \cong \big(\gr(\V^k(\g)[e^{-1}])\big)^{\J_\infty(\kfk \ltimes \g_{\geqslant 1})}. \]
	The first part of the statement is proved.
	
	We immediately deduce the isomorphisms 
	\[ \begin{aligned}
		\gr\big((\V^k(\g)[e^{-1}])^{\J_\infty(\kfk \ltimes \g_{\geqslant 1})}\big) & \cong \big(\J_\infty(\Sym \g)[e^{-1}]\big)^{\J_\infty(\kfk \ltimes \g_{\geqslant 1})}  \\
		& \cong  \big(\J_\infty(\Sym \g)[e^{-1}]\big)^{\J_\infty(K \ltimes G_{\geqslant 1})} \\
		& \cong \C[\g^\vee[e^{-1}] \git \J_\infty(K \ltimes G_{\geqslant 1})]
	\end{aligned} \]	
	by applying \Cref{lemma:coholology-pro-algebra}. As a consequence of \Cref{proposition:dilation-gan-ginzburg}, we have scheme isomorphisms
	\[ \J_\infty S_f \cong \g^\vee[e^{-1}] \git \J_\infty(K \ltimes G_{\geqslant 1}) \cong \J_\infty\big(\g^\vee[e^{-1}] \git (K \ltimes G_{\geqslant 1})\big).  \]
	This gives the second part of the statement.
\end{proof}

The following lemma gives a natural vertex algebra homomorphism between the invariant subalgebra $\big(\V^k(\g)[e^{-1}]\big)^{\J_\infty(\kfk \ltimes \g_{\geqslant 1})}$ and the chiral Hamiltonian reduction of the localisation $\V^k(\g)[e^{-1}]$. 

\begin{lemma} \label{lemma:coset-to-cohomology}
	The image of the graded vertex algebra homomorphism
	\[ \begin{aligned}
		(\V^k(\g)[e^{-1}])^{\J_\infty(\kfk \ltimes \g_{\geqslant 1})} & \longrightarrow \Cpx(\V^k(\g)[e^{-1}] = \V^k(\g)[e^{-1}] \otimes \F(\g_{\geqslant 1} \oplus (\g_{\geqslant 1})^\vee), \\
		v & \longmapsto v \otimes 1 
	\end{aligned} \]
	is in the kernel of the differential $\widehat\mu_{e^{-1}}(Q)_{(0)}$.
\end{lemma}

\begin{proof}
	Let $v$ be in $\V^k(\g)[e^{-1}]$. By using the explicit formula for $Q$ and Borcherds identity \eqref{equation:borcherds}, we get:
	\[ \widehat\mu_{e^{-1}}(Q)_{(0)} (v \otimes 1) = \sum_{i=1}^{d-1} (a_i \otimes \phi(a_i^\vee))_{(0)} (v \otimes 1) \]
	because $v$ commutes with the generators of $\F(\g_{\geqslant 1} \oplus (\g_{\geqslant 1})^\vee)$. Then we use the Borcherds identity \eqref{equation:borcherds}:
	\[ \begin{multlined}
		(a_i \otimes \phi(a_i^\vee))_{(0)} v = \sum_{n \geqslant 0} (a_i \otimes 1)_{(-n - 1)} \big((1 \otimes \phi(a_i^\vee))_{(n)} (v \otimes 1)\big) \\
		+ \sum_{n \geqslant 0} (1 \otimes \phi(a_i^\vee))_{(-n - 1)} \big((a_i{}_{(n)} v) \otimes 1\big).
	\end{multlined}  \] 
	In the previous formula, the first sum vanishes because $v$ commutes with the generators of $\F(\g_{\geqslant 1} \oplus (\g_{\geqslant 1})^\vee)$. By using the equality $a_i{}_{(n)} v = a_i t^n \cdot v$, we finally get:
	\[ \widehat\mu_{e^{-1}}(Q)_{(0)} (v \otimes 1) =  \sum_{n \geqslant 0} (1 \otimes \phi(a_i^\vee))_{(-n - 1)} \big((a_i t^n \cdot v) \otimes 1\big). \]
	By definition of being invariant, this quantity is zero if $v$ belongs to the invariant subalgebra $(\V^k(\g)[e^{-1}])^{\J_\infty(\kfk \ltimes \g_{\geqslant 1})}$.
\end{proof}

Now we can prove the main theorem.


\begin{proof}[Proof of \Cref{theorem:affine-w-algebra-coset}]
	Let us show that the localisation at $e$ does not modify the cohomology. By composing $\widehat{\mu}$ with the localisation homomorphism, we get another chiral comoment map
	\[ \widehat\mu_{e^{-1}} \colon \V(\g_{\geqslant 1}) \longrightarrow \V^k(\g)[e^{-1}]. \]
	Because the BRST complex construction is functorial, there is a natural cochain complex and vertex homomorphism
	\[ \Cpx(\V^k(\g)) \longrightarrow \Cpx(\V^k(\g)[e^{-1}]) = \Cpx(\V^k(\g))[e^{-1}] \]
	that is also a localisation homomorphism because tensoring with the fermionic vertex superalgebra commutes with localisation. Then we get a vertex superalgebra homomorphism by taking cohomologies:
	\[ \W^k(\g, f) = \Hgy\big(\Cpx(\V^k(\g))\big) \longrightarrow  \Hgy\big(\Cpx(\V^k(\g)[e^{-1}])\big). \]
	Let us prove that it is a vertex algebra isomorphism by using \Cref{proposition:reduction-quantisation}.

	It is straightforward to check that the moment maps $\mu\colon \g^\vee \rightarrow (\g_{\geqslant 1})^\vee$ and $\mu_{e^{-1}}\colon \g^\vee[e^{-1}] \rightarrow (\g_{\geqslant 1})^\vee$ are smooth by checking that their differential are surjective at any point. Moreover, $\J_\infty \mu^*$ coincides with $\gr \widehat{\mu}$ up to the isomorphism $\J_\infty(\Sym \g) \cong \gr \V^k(\g)$. \Cref{theorem:gan-ginzburg-isomorphism} and \Cref{proposition:hamiltonian-reduction-open} give the other assumptions that are necessary to apply \Cref{proposition:reduction-quantisation}. Therefore, both cohomologies of $\Cpx(\V^k(\g))$ and $\Cpx(\V^k(\g))[e^{-1}]$ are concentrated in degree $0$ and we have the following commutative diagram of vertex Poisson algebra homomorphisms:
	\[ \begin{tikzcd}
		\gr \Hgy^0(\Cpx(\V^k(g))) \arrow[r] \arrow[d, "\sim"'{sloped}] & \gr \Hgy^0\big(\Cpx(\V^k(\g))[e^{-1}]\big) \arrow[d, "\sim"{sloped}] \\
		\J_\infty \C[\mu^{-1}(0) \git G_{\geqslant 1}] \arrow[r, "\sim"] & \J_\infty \C[(\mu_{e^{-1}})^{-1}(0) \git G_{\geqslant 1}],
	\end{tikzcd} \]
	where the vertical arrows are isomorphisms. In fact, the lower horizontal arrow is an isomorphism by applying the arc functor $\J_\infty$ to the isomorphism from \Cref{proposition:hamiltonian-reduction-open}. Therefore, we get a vertex Poisson algebra isomorphism
	\[ \gr \W^k(\g, f) \cong \gr \Hgy^0\big(\Cpx(\V^k(\g)[e^{-1}]\big) \]
	that is given by the upper arrow. Since the filtrations on the cohomology vertex algebras are induced by the Li filtrations on the complexes, these filtrations are bounded from below and exhaustive. By \Cref{lemma:quantisation-iso}, we get the vertex algebra isomorphism:
	\[ \W^k(\g, f) \cong \Hgy^0\big(\Cpx(\V^k(\g)[e^{-1}]\big) = \Hgy\big(\Cpx(\V^k(\g)[e^{-1}]\big). \]
	
	As a consequence of \Cref{lemma:coset-to-cohomology}, we have a vertex algebra homomorphism
	\[ \widehat\lambda \colon (\V^k(\g)[e^{-1}])^{\J_\infty(\kfk \ltimes \g_{\geqslant 1})} \longrightarrow \Hgy^0\big(\Cpx(\V^k(\g)[e^{-1}]\big) \cong \W^k(\g, f). \]
	Both sides are equipped with exhaustive filtrations that are bounded from below, so it is an isomorphism if $\gr \widehat\lambda$ is an isomorphism, according to \Cref{lemma:quantisation-iso}. 
	
	By \Cref{proposition:gr-invariant-borisov-localisation} and \Cref{proposition:reduction-quantisation} (see the first part of this proof), we have the algebra isomorphisms
	\[ \begin{multlined}
		\gr \big(\big(\V^k(\g)[e^{-1}]\big)^{\J_\infty(\kfk \ltimes \g_{\geqslant 1})}\big) \cong \J_\infty \C[\g^\vee[e^{-1}] \git (K \ltimes G_{\geqslant 1})] \\ \text{and} \quad \gr \W^k(\g, f) \cong \J_\infty \C[\mu^{-1}(0) \git G_{\geqslant 1}]. 
	\end{multlined}\]
	Under these isomorphisms, $\gr \widehat\lambda$ is conjugate to the isomorphism
	\[ \J_\infty \C[\g^\vee[e^{-1}] \git (K \ltimes G_{\geqslant 1})] \cong  \J_\infty \C[\mu^{-1}(0) \git G_{\geqslant 1}] \]
	given by \Cref{theorem:slodowy-git}. Therefore $\gr \widehat\lambda$, and then $\widehat\lambda$, are isomorphisms.
\end{proof}



\subsection{Decomposition theorem for affine W-algebras}

\begin{lemma}
	The two vertex subalgebras 
	\[ \V^k(\kfk \ltimes \g_{\geqslant 1})[e^{-1}] \quad \text{and} \quad \big(\V^k(\g)[e^{-1}]\big)^{\J_\infty(\kfk \ltimes \g_{\geqslant 1})} \cong \W^k(\g, f) \]
	of $\V^k(\g)[e^{-1}]$ commute with each other.
\end{lemma}

\begin{proof}
	The lemma follows immediately from the following equalities:
	\[ \begin{aligned}
		& \big(\V^k(\g)[e^{-1}] \big)^{\J_\infty(\kfk \ltimes \g_{\geqslant 1})} \\
		= & \big\{ F \in \V^k(\g)[e^{-1}] \quad | \quad \text{for all} ~  x \in \kfk \ltimes \g_{\geqslant 1} ~\text{and} ~n \geqslant 0,  ~   xt^n \cdot F = x_{(n)} F = 0 \big\} \\
		= & \big\{ F \in \V^k(\g)[e^{-1}] \quad | \quad \text{for all} ~  x \in \V^k(\kfk \ltimes \g_{\geqslant 1})[e^{-1}],  ~   [x(z_1), F(z_2)] = 0 \big\}. 
	\end{aligned} \]
	The second equality follows because commuting with $\V^k(\kfk \ltimes \g_{\geqslant 1})[e^{-1}]$ is equivalent to commuting with $\kfk \ltimes \g_{\geqslant 1}$, since this space generates the non-localised algebra, and commuting with $e(z_1)$ implies commuting with $e^{-1}(z_1) = (e(z_1))^{-1}$ by \cite[Formula (3.1)]{arakawa2026faithful}.
	
\end{proof}

Consequently, there is a natural vertex algebra homomorphism
\[ \widehat{\delta} \colon \V^k(\kfk \ltimes \g_{\geqslant 1})[e^{-1}] \otimes \big(\V^k(\g)[e^{-1}]\big)^{\J_\infty(\kfk \ltimes \g_{\geqslant 1})} \longrightarrow \V^k(\g)[e^{-1}] \]
given by the $(-1)$-product. This provides the decomposition isomorphism for the affine W-algebra.


\begin{theorem} \label{theorem:decomposition-chiral}
	The vertex algebra homomorphism
	\[ \widehat{\delta} \colon \V^k(\kfk \ltimes \g_{\geqslant 1})[e^{-1}] \otimes \big(\V^k(\g)[e^{-1}]\big)^{\J_\infty(\kfk \ltimes \g_{\geqslant 1})} \longrightarrow \V^k(\g)[e^{-1}] \]
	is an isomorphism.
\end{theorem}

\begin{proof}	
	Both sides are equipped with increasing Li filtrations. Because these filtrations on both sides are bounded from below and exhaustive, we just need to prove that
	\[ \gr \widehat\delta \colon \gr \big(\V^k(\kfk \ltimes \g_{\geqslant 1})[e^{-1}]\big) \otimes \gr\big(\big(\V^k(\g)[e^{-1}]\big)^{\J_\infty(\kfk \ltimes \g_{\geqslant 1})}\big) \longrightarrow \gr \big(\V^k(\g)[e^{-1}]\big) \]
	is an isomorphism to apply \Cref{lemma:quantisation-iso}.
	
	By applying \Cref{proposition:gr-borisov-localisation}, we obtain the vertex Poisson algebra isomorphisms:
	\[ \gr\big(\V^k(\kfk \ltimes \g_{\geqslant 1})[e^{-1}]\big) \cong \J_\infty\Sym(\kfk \ltimes \g_{\geqslant 1})[e^{-1}]  \]
	and 
	\[ \gr\big(\V^k(\g)[e^{-1}]\big) \cong \J_\infty \Sym(\g)[e^{-1}]. \]
	
	Moreover, \Cref{proposition:gr-invariant-borisov-localisation} gives the vertex Poisson algebra isomorphism
	\[ \gr\big((\V^k(\g)[e^{-1}])^{\J_\infty(\kfk \ltimes \g_{\geqslant 1})}\big) \cong \J_\infty\C[\g^\vee[e^{-1}] \git(K \ltimes G_{\geqslant 1})].  \]
	
	Under these isomorphisms, $\gr \widehat{\delta}$ is conjugate to the vertex Poisson algebra isomorphism 
	\[ \J_\infty\delta^* \colon \J_\infty \C[(\kfk \ltimes \g_{\geqslant 1})^\vee][e^{-1}] \otimes \J_\infty\C[\g^\vee[e^{-1}] \git (K \ltimes G_{\geqslant 1})]  \overset{\sim}{\longrightarrow} \J_\infty \C[\g^\vee[e^{-1}]]  \]
	provided by \Cref{theorem:decomposition-slodowy}. This concludes the proof.
\end{proof}

\subsection{Inverse Hamiltonian reduction for the minimal W-algebra}

If we compose the localisation embedding $\V^k(\g) \hookrightarrow \V^k(\g)[e^{-1}]$ with the inverse $\widehat{\delta}^{-1}$ of the decomposition isomorphism (\Cref{theorem:decomposition-chiral}), we get an embedding of graded vertex algebras:
\[ \V^k(\g) \hooklongrightarrow \V^k(\kfk \ltimes \g_{\geqslant 1}) \otimes \W^k(\g, f). \]
This is similar to the inverse reduction embedding \eqref{equation:inverse-reduction}. To check whether it is an inverse reduction, we need to determine if the factor $\V^k(\kfk \ltimes \g_{\geqslant 1})$ is a \emph{localised Weyl vertex algebra}. 

We recall that the \emph{Weyl vertex algebra} of rank $r$ (also called \emph{$\beta\gamma$-system} in the physics literature) is a vertex algebra freely generated by $2 r$ generators $\beta_i$ and $\gamma_i$, where $1 \leqslant i \leqslant r$, whose only non-trivial OPE are given by:
\[ \beta_i(z_1) \gamma_i(z_2) \sim \frac{1}{z_1 - z_2}, \quad \text{for} \quad 1 \leqslant i \leqslant r. \]
We denote this vertex algebra by $\mathrm{Weyl}_r$. It can be equipped with a vertex algebra grading by imposing that the $\beta_i$ are in degree $1$ and the $\gamma_i$ are in degree $0$. In particular, the localisation $\mathrm{Weyl}_r[(\gamma_1)^{-1}]$ is well-defined. See for example \cite[Section 1.8]{desole2006finite}, where it corresponds to a purely even fermionic vertex algebra associated with the canonical symplectic form on $\C^r \times \C^r$. 

For $2 \leqslant i \leqslant d - 1$, we introduce the notations
\[ a_i \coloneqq x_{i + 1, 2} \quad \text{and} \quad b_i \coloneqq x_{1, i + 1}. \]
Then $\kfk \ltimes \g_{\geqslant 1}$ is spanned by $H$, $e$, $a_i$ and $b_i$ for $2 \leqslant i \leqslant d - 1$, with the only non-trivial bracket relations given by:
\[ [H, e] = e, \quad [H, b_i] = b_i, \quad [b_i, a_i] =e. \]

\begin{proposition} \label{proposition:relations-dual-w-algebra}
	The vertex algebra $\V^k(\kfk \ltimes \g_{\geqslant 1})[e^{-1}]$ is strongly generated by $H$, $e$, $e^{-1}$, $a_i$ and $b_i$ for $2 \leqslant i \leqslant d - 1$, with the only non-trivial OPE given by:
	\[ \begin{aligned}
		H(z_1) H(z_2) & \sim \frac{1}{(z_1 - z_2)^2} \frac{d-1}{d} k \id, \\
		H(z_1) e(z_2) & \sim \frac{1}{z_1 - z_2} e(z_2), \\
		H(z_1) e^{-1}(z_2) & \sim \frac{1}{z_1 - z_2} \big(- e(z_2)\big), \\
		H(z_1) b_i(z_2) & \sim \frac{1}{(z_1 - z_2)^1} b_i(z_2), \\
		b_i(z_1) a_i(z_2) & \sim \frac{1}{z_1 - z_2} e(z_2).
	\end{aligned}\]
\end{proposition}

\begin{proof}
	The only relation that must be computed is between $H$ and $e^{-1}$. This can be done with computations analogous to the proof of \Cref{proposition:relations-localisation-sl2}.
\end{proof}

From this, we deduce that this vertex algebra is a localised Weyl vertex algebra.

\begin{theorem} \label{theorem:commutant-localisation-weyl}
	The assignment 
	\[ \begin{aligned}
		\mathrm{Weyl}_{d-1} & \longrightarrow  \V^k(\kfk \ltimes \g_{\geqslant 1})[e^{-1}], \\
		\beta_1 & \longmapsto  \NO{e^{-1} H} + (3 - \frac{d - 1}{4 d} k) e^{-1} e^{-1} \partial e, \\
		\gamma_1 & \longmapsto e, \\
		\beta_i & \longmapsto - a_i, \\
		\gamma_i & \longmapsto e^{-1} b_i \quad \text{for} \quad 2 \leqslant i \leqslant d - 1,
	\end{aligned}  \]
	is a well-defined graded vertex algebra homomorphism which induces an isomorphism
	\[ \mathrm{Weyl}_{d-1}[(\gamma_1)^{-1}] \cong \V^k(\kfk \ltimes \g_{\geqslant 1})[e^{-1}]. \]
\end{theorem}

\begin{proof}
	The mapping is clearly degree preserving because $H$, the $a_i$ and $\partial e$ are of degree $1$, whereas $e$, $e^{-1}$ and the $b_i$ are of degree $0$.
	
	By using the Borcherds identity \eqref{equation:borcherds}, we get:
	\begin{equation} \label{equation:ope-quasi-beta-gamma} \begin{aligned}
		\NO{e^{-1} H}(z_1) e(z_2) & \sim \frac{1}{z_1 - z_2} \id, \\ 
		\NO{e^{-1} H}(z_1) (e^{-1} b_i)(z_2) & \sim 0, \\
		(e^{-1} b_i) (z_1) a_i(z_2) & \sim \frac{1}{z_1 - z_2} \id.
	\end{aligned} \end{equation}
	
	However, the OPE of $\NO{e^{-1}H}$ with itself is non-trivial:
	\[ \begin{multlined}
		\NO{e^{-1}H}(z_1) \NO{e^{-1}H}(z_2) \sim \frac{1}{z_1 - z_2} \big(12 - \frac{d-1}{d} k \big) (e^{-1} e^{-1} e^{-1} \partial e)(z_2) \\
		 + \frac{1}{(z_1 - z_2)^2} \big(\frac{d-1}{d} k - 12\big) e^{-1} e^{-1}.
	\end{multlined} \]
	By replacing $\NO{e^{-1}H}$ by the corrected term 
	\[ c \coloneqq  \NO{e^{-1}H} + (3 - \frac{d - 1}{4 d} k) e^{-1} e^{-1} \partial e, \]
	we get the same OPE as \eqref{equation:ope-quasi-beta-gamma} because $e$ and $e^{-1}$ commute with the $b_i$ and $a_i$. In addition, we get a trivial relation 
	\[ c(z_1) c(z_2) \sim 0. \]
	
	We have checked that the elements $c$, $e$, $e^{-1} b_i$ and $a_i$ have the OPE of the Weyl vertex algebra generators $\gamma_j$ and $\beta_j$. Therefore, we have a well-defined vertex algebra homomorphism.  
	
	To prove that the induced vertex algebra homomorphism \[ \mathrm{Weyl}_{d-1}[(\gamma_1)^{-1}] \longrightarrow \V^k(\kfk \ltimes \g_{\geqslant 1})[e^{-1}], \]
	is invertible, we can use the increasing Li filtrations on both sides. The associated graded algebra of $\mathrm{Weyl}_{d-1}$ is obtained by applying the arc functor $\J_\infty$ to a polynomial algebra with $2 (d - 1)$ generators denoted by $\beta_i$ and $\gamma_i$ (the algebra of polynomial functions over a symplectic vector space $\C^{d - 1} \times \C^{d - 1}$):
	\[ \gr \mathrm{Weyl}_{d-1} \cong \J_\infty \C[\gamma_i, \beta_i]_{1 \leqslant i \leqslant d - 1}. \]
	After taking $\gr$, the induced homomorphism is given by
	\[ \begin{aligned}
		\gr \mathrm{Weyl}_{d-1}[(\gamma_1)^{-1}] & \longrightarrow  \gr \V^k(\kfk \ltimes \g_{\geqslant 1})[e^{-1}], \\
		\beta_1 & \longmapsto  e^{-1} H, \\
		\gamma_1 & \longmapsto e, \\
		\beta_i & \longmapsto - a_i, \\
		\gamma_i & \longmapsto e^{-1} b_i \quad \text{for} \quad 2 \leqslant i \leqslant d - 1.
	\end{aligned}  \] 
	It is clearly an isomorphism whose inverse map is given by:
	\[  H \longmapsto \beta_1 \gamma_1, \quad e \longmapsto \gamma_1, \quad a_i \longmapsto - \beta_i, \quad b_i \longmapsto \gamma_1 \gamma_i.  \]
	Therefore, \Cref{lemma:quantisation-iso} implies that the vertex algebra homomorphism we have defined is an isomorphism.
\end{proof}

Therefore, we retrieve the inverse reduction for the minimal W-algebra in type~A as a corollary of \Cref{theorem:decomposition-chiral}.

\begin{theorem}[\cite{semikhatov1994inverting,adamovic2019realizations,adamovic2021realisation,fehily2024inverse}]
	There is a graded vertex algebra embedding
	\[ \V^k(\g) \hooklongrightarrow \mathrm{Weyl}_{d-1}[(\gamma_1)^{-1}] \otimes \W^k(\g, f). \]
\end{theorem}



